% Fibonacci-Pell nearest gaps: an all-exponent classification and
% quadratic-unit orbit rigidity
%
% Authors:
%   Dr. Denys Dutykh (Mathematics Department, Khalifa University of Science
%   and Technology, Abu Dhabi, UAE)
%   Prof. Laurent Vuillon (Univ. Savoie Mont Blanc, CNRS, LAMA, Chambery,
%   France)
%
\section{Effective classification of the even branch}
\label{sec:classification}

We now bound $q$, reduce both exponents by exact rational intervals, and
finish the even-exponent branch of \cref{thm:classification-intro}.
Until the terminal enumeration, we retain the standing hypothesis that an
even-exponent target exists with $q\geqslant7$.

\subsection{A moving-coefficient logarithmic form}

Quadratic conjugation of \cref{eq:negative-target}, followed by elimination
of $g$, gives the exact trace identity

\begin{equation}
 \lambda^n+\lambda^{-n-2h}
 =D_{q,\sigma}^2+\overline{D}_{q,\sigma}^{\,2}\lambda^{-2h}.
 \label{eq:trace-identity}
\end{equation}

Retain $\epsilon=(-1)^q$ and put

\begin{equation}
 \begin{aligned}
 u_\sigma&=\sigma+\sqrt2\varphi,
 &u'_\sigma&=-\sigma+\sqrt2\varphi^{-1},\\
 v_\sigma&=\sigma-\sqrt2\varphi,
 &v'_\sigma&=-\sigma-\sqrt2\varphi^{-1}.
 \end{aligned}
 \label{eq:four-coefficients}
\end{equation}

Binet's formula becomes

\begin{equation}
 \begin{aligned}
 D_{q,\sigma}
 &=\frac{u_\sigma\varphi^q+\epsilon u'_\sigma\varphi^{-q}}{\sqrt5},\\
 \overline D_{q,\sigma}
 &=\frac{v_\sigma\varphi^q+\epsilon v'_\sigma\varphi^{-q}}{\sqrt5}.
 \end{aligned}
 \label{eq:two-binet}
\end{equation}

The key is to retain both growing $\varphi^{2q}$ terms inside one
coefficient:

\begin{equation}
 \mathcal A_{\sigma,h}
 =\frac{u_\sigma^2+v_\sigma^2\lambda^{-2h}}5.
 \label{eq:moving-coefficient}
\end{equation}

Substitution in \cref{eq:trace-identity} gives

\begin{equation}
 \lambda^n=\mathcal A_{\sigma,h}\varphi^{2q}
 +\mathcal R_{\sigma,q,h},
 \label{eq:moving-main}
\end{equation}

where

\begin{align}
 \mathcal R_{\sigma,q,h}={}&
 \frac{2\epsilon}{5}
 \left((1-\sigma\sqrt2)
 +(1+\sigma\sqrt2)\lambda^{-2h}\right)
 \notag\\
 &+\frac{\varphi^{-2q}}5
 \left((u'_\sigma)^2+(v'_\sigma)^2\lambda^{-2h}\right)
 -\lambda^{-n-2h}.
 \label{eq:moving-remainder}
\end{align}

Here we used
$u_\sigma u'_\sigma=1-\sigma\sqrt2$ and
$v_\sigma v'_\sigma=1+\sigma\sqrt2$. Direct coefficient estimates give

\begin{equation}
 \abs{\mathcal R_{\sigma,q,h}}<4,
 \qquad
 \frac15<\mathcal A_{\sigma,h}<7.
 \label{eq:A-R-bounds}
\end{equation}

Since $\varphi^{14}=377\varphi+233>1597/2$, normalising
\cref{eq:moving-main} and using \cref{eq:A-R-bounds} yields the logarithmic
form

\begin{equation}
 \Lambda_1=n\log\lambda-2q\log\varphi
 -\log\mathcal A_{\sigma,h}
 \label{eq:Lambda1}
\end{equation}

with

\begin{equation}
 0<\abs{\Lambda_1}<21\varphi^{-2q}.
 \label{eq:Lambda1-small}
\end{equation}

The strict nonvanishing and the height of the moving coefficient must be
controlled uniformly. From
$\lambda^h=C_h+P_h\sqrt2$ one obtains

\begin{equation}
 \mathcal A_{\sigma,h}
 =\lambda^{-h}\frac{2\sqrt2}{5}
 \left((4+\sqrt5)P_h+2\sigma\varphi C_h\right).
 \label{eq:A-pell}
\end{equation}

Taking the full norm in $\Q(\sqrt2,\sqrt5)$ gives

\begin{equation}
 N_{K/\Q}(\mathcal A_{\sigma,h})
 =\frac{64}{625}
 \left(4+3P_h^2-2\sigma C_hP_h\right)^2.
 \label{eq:A-norm}
\end{equation}

If $\Lambda_1=0$, full norms would force $64I^2=625$ for an integer $I$,
which is impossible. We now record the height calculation. The inequalities
$h_{\mathrm{abs}}(\alpha\beta)\leqslant
h_{\mathrm{abs}}(\alpha)+h_{\mathrm{abs}}(\beta)$ and
$h_{\mathrm{abs}}(\alpha+\beta)\leqslant
h_{\mathrm{abs}}(\alpha)+h_{\mathrm{abs}}(\beta)+\log2$ give

\begin{equation}
 h_{\mathrm{abs}}(u_\sigma),\,
 h_{\mathrm{abs}}(v_\sigma)
 \leqslant
 \log2+\frac{\log2}{2}+\frac{\log\varphi}{2}
 <\frac32.
\end{equation}

Since $h_{\mathrm{abs}}(\lambda)=\tfrac12\log\lambda$, the definition
\cref{eq:moving-coefficient} yields

\begin{align}
 h_{\mathrm{abs}}(\mathcal A_{\sigma,h})
 &\leqslant
 2h_{\mathrm{abs}}(u_\sigma)
 +2h_{\mathrm{abs}}(v_\sigma)
 +2h\,h_{\mathrm{abs}}(\lambda)+\log10\notag\\
 &<h+9<h+11.
\end{align}

Together with \cref{eq:A-R-bounds}, this proves

\begin{equation}
 \height{\mathcal A_{\sigma,h}}<h+11,
 \qquad
 \abs{\log\mathcal A_{\sigma,h}}<2.
 \label{eq:A-height}
\end{equation}

Thus \cref{lem:matveev} applies to \cref{eq:Lambda1} with

\begin{equation}
 A_1=2\log\lambda,
 \qquad
 A_2=2\log\varphi,
 \qquad
 A_3=4h+44,
 \qquad
 B=2q.
\end{equation}

Writing again $H=1+\log(2q)$ and using
\cref{eq:matveev-prefactor}, comparison with
\cref{eq:Lambda1-small} gives

\begin{equation}
 2q\log\varphi
 <9.275\mathbin{\cdot}10^{12}(4h+44)H+\log21.
 \label{eq:q-h-comparison}
\end{equation}

By \cref{eq:h-log-bound} and $\log\varphi>2/5$,

\begin{equation}
 q<\frac54\left[
 9.275\mathbin{\cdot}10^{12}
 \left(9.6\mathbin{\cdot}10^{14}H+44\right)H+4
 \right].
 \label{eq:q-self-bound}
\end{equation}

At $q_0=10^{32}$, one has $H(q_0)<79$ and the right-hand side of
\cref{eq:q-self-bound} is less than $7\mathbin{\cdot}10^{31}$. After
division by $q$, every term on the right is decreasing for
$q\geqslant q_0$. Consequently

\begin{equation}
 q<10^{32}.
 \label{eq:q-1e32}
\end{equation}

The moving coefficient \cref{eq:moving-coefficient} is the central
mechanism of the proof. Its error decays with $q$, while its height grows
only linearly with $h$, which is already logarithmic in $q$.

\subsection{An absolute-value reduction lemma}

We use a symmetric form of the Baker--Davenport reduction, closely related
to the lemma of Dujella and Peth\H{o} \cite{DujellaPetho1998}. We include a
short proof to avoid a hidden sign condition.

\begin{lemma}[Absolute-value reduction]
\label{lem:absolute-reduction}
Let $\gamma,\mu\in\R$, let $p\in\Z$, let $Q\in\Z_{>0}$, and let
$M\geqslant1$. Suppose that $\abs{Q\gamma-p}<1/2$ and that integers
$a,b$ satisfy

\begin{equation}
 0<\abs{a\gamma-b+\mu}<AB^{-w},
 \qquad
 1\leqslant a\leqslant M,
 \label{eq:reduction-hypothesis}
\end{equation}

where $A>0$, $B>1$, and $w>0$ are real. Put

\begin{equation}
 \eta=\dist{Q\mu}-M\abs{Q\gamma-p},
 \label{eq:eta}
\end{equation}

where $\dist{x}$ is the distance from $x$ to the nearest integer. If
$\eta>0$, then

\begin{equation}
 w<\frac{\log(AQ/\eta)}{\log B}.
 \label{eq:reduction-conclusion}
\end{equation}
\end{lemma}

\begin{proof}
Write $\dist{x}$ for the distance to the nearest integer. Since
$ap-bQ\in\Z$, the triangle inequality on $\R/\Z$ gives

\begin{align}
 \dist{Q\mu}
 &\leqslant
 \dist{Q\mu+a(Q\gamma-p)}+a\abs{Q\gamma-p}\notag\\
 &\leqslant
 Q\abs{a\gamma-b+\mu}+M\abs{Q\gamma-p}.
\end{align}

The strict hypothesis therefore gives $\eta<AQB^{-w}$, which is equivalent
to \cref{eq:reduction-conclusion}.
\end{proof}

The point of this symmetric form is that no sign of the logarithmic form is
assumed. We apply it first to the fixed shift, reducing $h$ or $q$, and then
to the moving shift, which supplies the absolute bound for $q$.

\subsection{Certified reductions of the two exponents}

Set

\begin{equation}
 \gamma=\frac{\log\varphi}{\log\lambda},
 \qquad
 \mu_\sigma
 =\frac{\log(u_\sigma^2/5)}{\log\lambda}.
 \label{eq:gamma-mu}
\end{equation}

Put

\begin{equation}
 \delta_\sigma
 =\epsilon\frac{u'_\sigma}{u_\sigma}\varphi^{-2q},
 \qquad
 \widetilde\delta_\sigma
 =\epsilon\frac{v'_\sigma}{v_\sigma}\varphi^{-2q}.
\end{equation}

Dividing the Binet expansion of \cref{eq:trace-identity} by
$(u_\sigma^2/5)\varphi^{2q}$ gives the exact identity

\begin{align}
 e^{\Lambda_0}-1
 ={}&
 2\delta_\sigma+\delta_\sigma^2
 +\frac{v_\sigma^2}{u_\sigma^2}
  (1+\widetilde\delta_\sigma)^2\lambda^{-2h}\notag\\
 &-\frac5{u_\sigma^2}
  \lambda^{-n-2h}\varphi^{-2q}.
 \label{eq:fixed-exact-expansion}
\end{align}

For both signs, direct comparison gives

\begin{equation}
 1<u_\sigma<4,\qquad
 \abs{\frac{u'_\sigma}{u_\sigma}}<\frac32,\qquad
 \abs{\frac{v_\sigma}{u_\sigma}}<3,\qquad
 \abs{\frac{v'_\sigma}{v_\sigma}}<\frac32.
\end{equation}

Since $q\geqslant7$, $\varphi^{14}>1597/2$, so both delta terms have
absolute value below $1/100$. The three terms in
\cref{eq:fixed-exact-expansion} are therefore bounded, respectively, by
$4\varphi^{-2q}$, $10\lambda^{-2h}$, and $5\varphi^{-2q}$. Moreover,
$h\geqslant3$ and $\lambda^6>190$, whence

\begin{equation}
 \abs{e^{\Lambda_0}-1}
 <10\left(\lambda^{-2h}+\varphi^{-2q}\right)<\frac1{10}.
\end{equation}

The inequality $\abs{\log(1+x)}\leqslant\abs{x}/(1-\abs{x})$ for
$\abs{x}<1$ now gives the sharper estimate

\begin{equation}
 0<\abs{\Lambda_0}
 <12\left(\lambda^{-2h}+\varphi^{-2q}\right),
 \label{eq:Lambda0-sharp}
\end{equation}

where $\Lambda_0$ is defined in \cref{eq:Lambda0}. Put
$w=\min\{h,q\gamma\}$. Since $\log\lambda>4/5$,
\cref{eq:Lambda0-sharp} implies

\begin{equation}
 0<\abs{2q\gamma-n+\mu_\sigma}
 <30(\lambda^2)^{-w}.
 \label{eq:fixed-reduction-form}
\end{equation}

Take

\begin{equation}
 M_0=2\mathbin{\cdot}10^{32}
 \label{eq:M0}
\end{equation}

and the rational number

\begin{equation}
 \frac pQ
 =\frac{2585762774943540084566744409566833}
 {4736007914708481810221186278309584}.
 \label{eq:common-convergent}
\end{equation}

Exact rational intervals for square roots and logarithms certify

\begin{equation}
 \gcd(p,Q)=1,
 \qquad
 Q>6M_0,
 \qquad
 2Q\abs{Q\gamma-p}<1.
 \label{eq:convergent-certificate}
\end{equation}

Only the certified error is used below. For both signs, the same interval
certificate gives

\begin{equation}
 \eta_\sigma
 =\dist{Q\mu_\sigma}-M_0\abs{Q\gamma-p}>\frac7{100},
 \label{eq:fixed-eta}
\end{equation}

and

\begin{equation}
 \eta_\sigma(\lambda^2)^{48}>30Q.
 \label{eq:fixed-cutoff}
\end{equation}

Apply \cref{lem:absolute-reduction} with

\begin{equation}
 a=2q,\quad b=n,\quad \mu=\mu_\sigma,\quad
 A=30,\quad B=\lambda^2,\quad
 w=\min\{h,q\gamma\},\quad M=M_0.
\end{equation}

\Cref{eq:convergent-certificate,eq:fixed-eta,eq:fixed-cutoff}
then prove $\min\{h,q\gamma\}<48$. Since $\gamma>1/2$, either $h<48$,
or $q<96$ and then

\begin{equation}
 h<n<2q<192.
\end{equation}

Thus every target with $q\geqslant7$ satisfies

\begin{equation}
 3\leqslant h<192,
 \qquad
 h\ \text{odd}.
 \label{eq:h-192}
\end{equation}

For each of the $190$ pairs

\begin{equation}
 \sigma\in\{1,-1\},
 \qquad
 h\in\{3,5,\ldots,191\},
\end{equation}

put

\begin{equation}
 \nu_{\sigma,h}
 =\frac{\log\mathcal A_{\sigma,h}}{\log\lambda}.
\end{equation}

\Cref{eq:Lambda1-small} gives

\begin{equation}
 0<\abs{2q\gamma-n+\nu_{\sigma,h}}
 <27(\varphi^2)^{-q}.
 \label{eq:moving-reduction-form}
\end{equation}

The same $p/Q$ satisfies, for every one of the $190$ shifts,

\begin{equation}
 \eta_{\sigma,h}
 =\dist{Q\nu_{\sigma,h}}-M_0\abs{Q\gamma-p}>\frac9{1000}
 \label{eq:moving-eta}
\end{equation}

and

\begin{equation}
 \eta_{\sigma,h}(\varphi^2)^{90}>27Q.
 \label{eq:moving-cutoff}
\end{equation}

The weakest certified interval occurs at $(\sigma,h)=(1,31)$ and remains
strictly above the displayed margin. Apply \cref{lem:absolute-reduction}
for each pair with

\begin{equation}
 a=2q,\quad b=n,\quad \mu=\nu_{\sigma,h},\quad
 A=27,\quad B=\varphi^2,\quad w=q,\quad M=M_0.
\end{equation}

\Cref{eq:convergent-certificate,eq:moving-eta,eq:moving-cutoff}
prove

\begin{equation}
 q<90.
 \label{eq:q-90}
\end{equation}

The interval proof uses only rational arithmetic. For a positive rational
$x$, logarithms are enclosed using

\begin{equation}
 \log x
 =2\sum_{r=0}^{N-1}\frac{z^{2r+1}}{2r+1}+R_N,
 \qquad
 z=\frac{x-1}{x+1},
 \label{eq:log-series}
\end{equation}

with the geometric tail bound

\begin{equation}
 \abs{R_N}
 \leqslant\frac{2\abs z^{2N+1}}{(2N+1)(1-z^2)}.
 \label{eq:log-tail}
\end{equation}

Square roots are bounded between consecutive rational grid points. Hence no
binary floating-point comparison enters
\cref{eq:convergent-certificate,eq:fixed-eta,eq:fixed-cutoff,eq:moving-eta,eq:moving-cutoff}.

At this point the unbounded even-exponent problem has become finite: every
hypothetical target has $q<90$. What remains is an exact recurrence
calculation, not a further asymptotic estimate.

\subsection{Exact terminal enumeration}

The remaining lower boundary $q=1$ is direct:

\begin{equation}
 D_{1,1}=\lambda,
 \qquad
 D_{1,-1}=\lambda^{-1}.
 \label{eq:q1-boundary}
\end{equation}

For the plus sign, the first integral power strictly above
$D_{1,1}^2=\lambda^2$ is $\lambda^3>2\lambda^2$. For the minus sign, every
positive integral power is at least
$\lambda>2\lambda^{-2}=2D_{1,-1}^2$. Thus neither sign has a strict
candidate at $q=1$.

It remains to test $2\leqslant q<90$, $3\nmid q$, and both signs. For each
row, exact recurrence arithmetic constructs $D_{q,\sigma}^2$ and the first
even Pell unit above it. The algorithm then checks the strict upper window,
computes the two coefficient gcd, and tests the exact norm identity. Signs of
$a+b\sqrt2$ are decided over the integers from $a$, $b$, and
$a^2-2b^2$; no numerical approximation is used.

The complete output consists of

\begin{align}
 q=2, \sigma=-1:
 &\quad
 \lambda^2-(-F_2+F_3\sqrt2)^2
 =6\lambda^{-1},
 \label{eq:low-q2}\\
 q=4, \sigma=1:
 &\quad
 \lambda^6-(F_4+F_5\sqrt2)^2
 =40\lambda.
 \label{eq:low-q4}
\end{align}

No other row passes. In particular, the range $7\leqslant q<90$ is empty;
together with \cref{eq:q-90}, this proves the even-exponent branch of
\cref{thm:classification-intro}. The odd branch cannot use the primitive
rectangle or the inequality $h<n$; the next section closes it by reversing
the order of the fixed and moving logarithmic-form estimates.
